\section{Evaluation Results}
\label{sec:results}

We report the results by question rather than by study: which decisions a controller should take from the model (\S\ref{sec:rq-stop}), what remains for the model once the controller applies the verifier's evidence standard (\S\ref{sec:fair}), where the resulting security boundary lies (\S\ref{sec:adversarial}), where the controller's knowledge comes from and how it fails when that knowledge is wrong (\S\ref{sec:mismatch}), and whether the effects replicate on the second family (\S\ref{sec:external}). Appendix~\ref{app:tables} gives every cell of every study.

\subsection{Probing and Stopping Are Separate Decisions}
\label{sec:rq-stop}
\label{sec:rq-rank}
A controller that picks probes changes two things at once: it removes the model's own choice, and it replaces that choice with a ranking. The decomposition study separates the two with blind configurations, in which the planner receives the same prompt whatever the selector; the stopping study then crosses who selects probes with who may stop. \Cref{fig:probe-stop} shows the second.

\begin{figure*}[t]
  \centering
  \includegraphics[width=0.92\textwidth]{fig-probe-stop.pdf}
  \caption{Stopping study: verified completion by who selects probes and who may stop (72 episodes per point). The dashed line is the post hoc LLM-free controller (EIG/cost selection and controller stop, same seeds and tasks).}
  \label{fig:probe-stop}
\end{figure*}

First, the ranking itself helps. With identical planner information, EIG/cost selection beats random selection by $+0.306$ $[+0.14,+0.47]$ on Qwen2.5-7B (a pre-registered endpoint), and by $+0.26$ to $+0.35$ on every other model that concludes. Taking the choice away, in contrast, helps weak models and hurts strong ones: random selection instead of the model's own choices raises completion by $+0.556$ for Qwen2.5-7B and lowers it by 0.222 and 0.278 for Qwen2.5-32B and 72B. Showing the ranking to the planner changes completion by at most $+0.06$.

Second, selecting probes does not make a model conclude. Under our prompt and serving configuration, Llama-3.1-8B completed none of its 360 decomposition episodes: all 3{,}389 of its decisions request another probe, in valid JSON. With a controller stop it rises from 0 to 0.917 ($+0.917$ $[+0.79,+1.00]$, pre-registered), and controller selection then adds $+0.056$ $[-0.07,+0.18]$, no detectable completion, although it halves probe cost (2.35 against 4.78 units). Qwen2.5-7B needs both parts: a stop alone lifts it from 0.083 to 0.694, and selection adds $+0.222$ $[+0.08,+0.39]$. A controller cannot know in advance which part a model lacks, so it has to supply both.

Third, the posterior stop concludes by elimination. With controller selection and stop, four of the five models fail on the same six episodes: a correct DNS command-and-control verdict reached by eliminating the alternatives, which the verifier rejects for lack of a confirming signature. Strong models that decide when to stop, in configurations without the controller stop, keep probing until that signature appears. The LLM-free controller with the posterior stop completes 0.917 at 2.35 units, and four models score exactly the same when the controller both selects and stops.

\subsection{What Remains for the Model}
\label{sec:fair}
The comparison above mixes two things: whether an LLM helps, and which evidence standard ends the episode. The v1.1 study removes this confound. It gives the LLM-free controller the confirmation stop of \S\ref{sec:finish}, adds a fixed playbook and catalogue order as simpler references, and runs every LLM configuration with the finish guard, on both families with the same seeds, tasks, and budgets. Its primary endpoints compare Qwen2.5-72B deciding when to stop against the LLM-free confirmation stop (97.5\% intervals). \Cref{tab:v11fair} reports both episodes that meet the evidence policy and episodes that name the right cause.

\begin{table}[t]
  \caption{What remains for the model (v1.1): episodes that meet the evidence policy / episodes that name the right cause, when the LLM decides when to stop (``LLM''), under the posterior stop (``post.''), and under the confirmation stop (``conf.''). Every LLM row uses controller selection and the finish guard; seeds, tasks, and budgets match the LLM-free row.}
  \label{tab:v11fair}
  \centering\small
  \begin{adjustbox}{max width=\columnwidth}% generated by paper/make_tables.py -- do not edit
\begin{tabular}{llrrrr}
\toprule
 & & \multicolumn{2}{c}{sec-triage (72)} & \multicolumn{2}{c}{sigma-triage (152)} \\
\cmidrule(lr){3-4}\cmidrule(lr){5-6}
Planner & Controller & verified & cost & verified & cost \\
\midrule
LLM-free & posterior stop & 66 & 2.35 & 152 & 2.20 \\
LLM-free & confirmation stop & 72 & 2.54 & 152 & 2.20 \\
LLM-free & fixed playbook, conf. & 66 & 2.90 & 152 & 2.20 \\
LLM-free & catalogue order, conf. & 65 & 4.93 & 151 & 2.96 \\
\addlinespace
Q-7B & guard, LLM stops & 66 & 7.17 & 126 & 3.22 \\
Q-7B & guard + posterior stop & 66 & 2.35 & 141 & 2.12 \\
Q-7B & guard + confirmation stop & 72 & 2.54 & 141 & 2.12 \\
\addlinespace
L-8B & guard, LLM stops & 0 & 10.00 & 1 & 9.78 \\
L-8B & guard + posterior stop & 66 & 2.35 & 152 & 2.20 \\
L-8B & guard + confirmation stop & 72 & 2.54 & 152 & 2.20 \\
\bottomrule
\end{tabular}
\end{adjustbox}
\end{table}

First, with the verifier's standard, the controller alone completes every case: 72 of 72 sec-triage and 152 of 152 sigma-triage episodes, at 2.54 and 2.20 units, 0.19 units more than the posterior stop on sec-triage. Adaptive selection still matters on sec-triage, where a fixed playbook reaches 66 and catalogue order 65, but not on sigma-triage, where every order completes at least 151 of 152.

Second, no model adds to that controller. Qwen2.5-72B deciding when to stop meets the evidence policy in 72 of 72 sec-triage episodes ($+0.000$ $[0.000, 0.000]$) and 148 of 152 sigma-triage episodes ($-0.028$ $[-0.111, 0.000]$; the mean over rules of per-rule differences, $-4/152 = -0.026$ per episode), at 3.86 and 3.25 units against 2.54 and 2.20. By the interpretation written before the run, we measured no gain from the LLM in deciding when to stop. Qwen2.5-32B and Llama-3.1-70B behave alike, and every sigma-triage miss of the three large models names the right cause without a confirming signature. When the model decides when to stop, Llama-3.1-8B completes 0 and 1 episodes and Qwen2.5-7B 66 and 126. Qwen2.5-7B also falls short under the controller stop (141 of 152), for one reason: on three rules it proposes a finish at a posterior of 0.794 to 0.796, just below the guard's threshold, the guard rejects it, and the model repeats the proposal until the decision budget is spent (11 episodes, 9 to 11 rejections each). The controller stop cannot fire below 0.8, so a planner can stall a guarded controller. The v1.2 study adds a recovery rule, under which the controller runs its own top-ranked probe after two consecutive rejected finishes: a scripted planner that insists on finishing at a score of 0.5 stalls all 152 sigma-triage episodes without it and completes all 152 with it.

Third, the evidence policy is a trade-off between automation and safety, not a detail of scoring. Our verifier accepts a verdict only if one cited observation carries a signature that the cause alone produces, so a right cause reached by combining or excluding observations counts as a failure; \cref{tab:v11fair} shows that under a ``right cause'' standard the posterior stop would also complete 72 of 72. The v1.2 study measures the alternatives with the LLM-free controller (\cref{tab:v12policy}). It adds \emph{comp-triage}, a third family of three scenarios in which an attack and three legitimate activities use the same tools, outcomes are probabilistic, and five of twelve causes have no single-observation signature; a \emph{joint} verifier that accepts any combination of observations that, under the true distributions, favours the cause at least tenfold over every other named cause; and stop rules that require joint evidence against the named causes, or also against \textit{other}. On comp-triage, joint evidence names the right cause in 59 of 72 episodes against 20 for the single-signature rule, which escalates 52. The price is paid in the attacker's currency: joint evidence closed 4 of 18 attacks as the paired legitimate activity, and when the true cause is absent from the candidates it named another actionable cause in 24 of 48 sec-triage episodes. Requiring the evidence to beat \textit{other} tenfold did not change this, and thirtyfold only reduced it to 18. Only the single-signature rule named no wrong cause in any setting, at the cost of escalating most composite cases; lowering its ratio to 3 moves along the same frontier (37 right causes, 1 missed attack). A deployment has to choose a point on this frontier; ours favours escalation.

\begin{table*}[t]
  \caption{Evidence policies (v1.2, LLM-free controller, EIG/cost selection; counts). ``single'' / ``joint'': episodes accepted by the single-signature and the joint-identification verifier. sec, cause absent: the true cause is removed from the candidate set. sec, confusion: each cause's table rows mixed 0.25 toward one other cause. comp-triage: three scenarios pairing an attack with legitimate activities that use the same tools, probabilistic outcomes, five of twelve causes without a single-observation signature.}
  \label{tab:v12policy}
  \centering\small
  \begin{adjustbox}{max width=\textwidth}% generated by paper/make_tables.py -- do not edit
\begin{tabular}{lrrrrrrrrr}
\toprule
 & \multicolumn{2}{c}{sec, matched (48)} & \multicolumn{2}{c}{sec, cause absent (48)} & sec, confusion & \multicolumn{4}{c}{comp-triage (72; 18 attacks)} \\
\cmidrule(lr){2-3}\cmidrule(lr){4-5}\cmidrule(lr){6-6}\cmidrule(lr){7-10}
Stop rule & single & joint & wrong & escal. & verified & right cause & wrong & missed attack & escal. \\
\midrule
posterior (0.8) & 42 & 48 & 24 & 24 & 47 & 64 & 7 & 5 & 1 \\
single signature, $r{=}10$ & 48 & 48 & 0 & 48 & 0 & 20 & 0 & 0 & 52 \\
single signature, $r{=}3$ & 48 & 48 & 0 & 48 & 0 & 37 & 2 & 1 & 33 \\
joint, named causes, $r{=}10$ & 42 & 48 & 24 & 24 & 47 & 59 & 6 & 4 & 7 \\
joint, incl.\ \textit{other}, $r{=}10$ & 42 & 48 & 24 & 24 & 47 & 59 & 6 & 4 & 7 \\
joint, incl.\ \textit{other}, $r{=}30$ & 42 & 48 & 18 & 30 & 18 & 55 & 2 & 1 & 15 \\
\bottomrule
\end{tabular}
\end{adjustbox}
\end{table*}

Fourth, the stopping failure is partly a prompting effect. Replaying 60 archived requests with overwhelming evidence or no probe left, Llama-3.1-8B returned valid, untruncated JSON every time and asked for a probe in 60, 60, and 58 of 60 under the original prompt, a prompt with an example finish, and a prompt that states when to finish. In live episodes inside \hesp{} the explicit rule made it finish in 8 of 48, the other variants in none; for Qwen2.5-7B both variants raised completion inside \hesp{} from 43 to 48 of 48. We did not give the Memory-only and ReAct-style baselines an equally developed prompt, so the size of their gap is specific to the prompt we used.

\subsection{The Security Boundary}
\label{sec:adversarial}
\label{sec:security}
The threat model allows the attacker to control part of the evidence. Without adversarial evidence, the ledger alone already removes most wrong verdicts: in the table-source study, the ReAct-style Qwen2.5-32B and 72B baselines called 13 of 51 and 20 of 54 actionable episodes benign among those they ruled on, Memory-only 2 of 54 and 1 of 53, and \hesp{} none. The adversarial studies test an instruction to close the case as an authorized scan, carried in free-text fields (four texts that differ in wording and position), and forged evidence: a source-reputation probe that reports a known scanner (\emph{forged probe}) and an upstream reputation feed that two probes read (\emph{forged feed}). \Cref{tab:v11attack} attributes the outcomes to components on all five models; the v1.0 results are in \cref{tab:v10c} (Appendix~\ref{app:tables}).

\begin{table*}[t]
  \caption{The security boundary (v1.1, sec-triage, five models): episodes that close an actionable case as benign (``missed'') or end without a verdict (``unresolved'') under four injection texts (of 72 per model); episodes closed as benign under forged evidence (of 18; range over the five models); verified episodes on the two benign causes (of 6; range over the five models).}
  \label{tab:v11attack}
  \centering\small
  \begin{adjustbox}{max width=\textwidth}% generated by paper/make_tables.py -- do not edit
\begin{tabular}{lrrrrrrrr}
\toprule
 & \multicolumn{5}{c}{4 injection texts: missed / unresolved (of 72)} & \multicolumn{2}{c}{forged: missed (of 18)} & benign \\
\cmidrule(lr){2-6}\cmidrule(lr){7-8}\cmidrule(lr){9-9}
Configuration & Q-7B & Q-32B & Q-72B & L-8B & L-70B & probe & feed & verif.\ (of 6) \\
\midrule
Memory-only & 69 / 0 & 57 / 0 & 69 / 0 & 0 / 72 & 63 / 0 & 0--6 & 0--15 & 0--6 \\
Memory-only, raw text removed & 0 / 65 & 0 / 0 & 0 / 0 & 0 / 72 & 0 / 24 & 0--6 & 0--15 & 0--6 \\
\hesp{}, guard only & 0 / 44 & 0 / 46 & 0 / 66 & 0 / 72 & 0 / 58 & 0 & 0 & 0--6 \\
\hesp{}, stop, no guard & 39 / 0 & 36 / 0 & 48 / 0 & 0 / 0 & 36 / 0 & 18 & 18 & 6 \\
\hesp{}, guard + stop & 0 / 24 & 0 / 29 & 0 / 42 & 0 / 0 & 0 / 2 & 18 & 18 & 6 \\
\quad + corroboration by probe & 0 / 0 & 0 / 0 & 0 / 11 & 0 / 0 & 0 / 0 & 0 & 0 & 3 \\
\quad + corroboration by source & 0 / 12 & 0 / 0 & 0 / 3 & 0 / 0 & 0 / 0 & 0 & 0 & 3 \\
\bottomrule
\end{tabular}
\end{adjustbox}
\end{table*}

First, the finish guard is what blocks injected dismissals. Every model that decides and reaches its own verdict follows the instruction, and controller selection and stop without the guard still let the planner finish first with the injected verdict in 36 to 48 of 72 episodes. Every guarded configuration closed 0 of 72 on every model; in v1.0 the pre-registered comparison on Qwen2.5-7B was $+1.000$ $[+1.000,+1.000]$, all 18 paired episodes moving the same way. Removing the raw text from the prompt also blocks the dismissal, at no cost for Qwen2.5-32B and 72B but leaving 65 of 72 Qwen2.5-7B cases unresolved. The injection still costs analyst time: the guard without a controller stop leaves 44 to 72 of 72 episodes unresolved, and adding the stop 0 to 42.

Second, forged evidence defeats the controller. A single forged scanner answer lifts the authorized scan above 0.8 after one probe, and the controller stop concludes: 18 of 18 actionable cases are closed as benign on every model, and a forged feed read by two probes does the same. Acting on the most informative evidence first is exactly what the attacker exploits. The forged feed also misleads the larger LLMs themselves (up to 15 of 18 under Memory-only).

Third, corroboration for benign verdicts closes this gap at a price. Requiring two specifically supporting observations escalates all forged cases on every model, but the per-probe rule survived the forged feed only because the controller's second probe read a different source. A targeted v1.2 test makes a fixed playbook read \texttt{source\_ips} and then \texttt{threat\_intel}, both fed by the forged source: per-probe corroboration then closes 18 of 18 actionable cases as benign, and per-source-group corroboration escalates all 18. Both rules halve completion on the two honest benign causes of sec-triage (3 of 6), because the monitoring false positive has only one specific signature there; across the 26 benign explanations of sigma-triage, corroboration escalates 8 of 104 benign episodes. All counts come from one environment, and a 0 of 18 cell bounds attack success only loosely (about 15\%, one-sided 95\%, even if its six causes and three repeats were independent).

\subsection{Where the Knowledge Comes From}
\label{sec:rq-help}
\label{sec:mismatch}
The controller's knowledge sits in its prediction tables. Counted tables match the generator's: with $k{=}20$ development episodes per cause and variant, \hesp{} with the finish guard lifts Qwen2.5-7B from 0.125 to 1.000 ($+0.875$ $[+0.74,+0.97]$, pre-registered), equal to designer tables on all three Qwen models, and $k{=}5$ already suffices (\cref{tab:v06}, Appendix~\ref{app:tables}), which reflects how deterministic the environment is. Counted tables cost 1.2 to 2.4 more units, all of it in the row for \textit{other}, the one cause development data never exhibits. The models cannot write the tables themselves: tables elicited from the planner leave completion at 0.069 to 0.375.

In our environments the tables come from the same generator as the test episodes. \Cref{tab:v10b} breaks that link with the LLM-free controller. Three findings stand out; the first two hold on both families.

\begin{table*}[t]
  \caption{Table error (LLM-free, posterior stop at 0.8). Under EIG/cost selection: the shares of episodes that are verified, name the right cause without a confirming signature, name a wrong cause, and escalate (these sum to one); under random selection: verified completion. ``Perturbed'' mixes every table row with a random Dirichlet(1) row with weight $\lambda$ (20 random tables per level); ``one cause missing'' averages over removing each cause in turn from the development data and testing on that cause.}
  \label{tab:v10b}
  \centering\small
  \begin{adjustbox}{max width=\textwidth}% generated by paper/make_tables.py -- do not edit
\begin{tabular}{lrrrrrrrrrr}
\toprule
 & \multicolumn{5}{c}{sec-triage} & \multicolumn{5}{c}{sigma-triage} \\
\cmidrule(lr){2-6} \cmidrule(lr){7-11}
Tables & verif. & correct, unverif. & wrong & esc. & random verif. & verif. & correct, unverif. & wrong & esc. & random verif. \\
\midrule
matched tables & 0.917 & 0.083 & 0.000 & 0.000 & 0.792 & 1.000 & 0.000 & 0.000 & 0.000 & 0.886 \\
base variant only & 0.917 & 0.083 & 0.000 & 0.000 & 0.792 & 1.000 & 0.000 & 0.000 & 0.000 & 0.886 \\
$k{=}1$, base only & 0.917 & 0.083 & 0.000 & 0.000 & 0.778 & 1.000 & 0.000 & 0.000 & 0.000 & 0.882 \\
perturbed, $\lambda{=}0.25$ & 0.885 & 0.083 & 0.000 & 0.031 & 0.500 & 0.992 & 0.008 & 0.000 & 0.000 & 0.910 \\
perturbed, $\lambda{=}0.5$ & 0.494 & 0.177 & 0.000 & 0.329 & 0.227 & 0.978 & 0.017 & 0.000 & 0.005 & 0.838 \\
perturbed, $\lambda{=}0.75$ & 0.098 & 0.027 & 0.004 & 0.871 & 0.042 & 0.805 & 0.040 & 0.009 & 0.146 & 0.513 \\
one cause missing (mean) & 0.000 & 0.000 & 0.500 & 0.500 & 0.000 & 0.000 & 0.000 & 0.079 & 0.921 & 0.000 \\
\bottomrule
\end{tabular}
\end{adjustbox}
\end{table*}

First, imprecise tables make the controller fail by escalating. As $\lambda$ grows, completion falls, but the wrong-cause rate stays at or below 0.009 even when three quarters of every row is noise, and EIG/cost stays ahead of random selection at every level. These results concern random-mixture perturbations and do not establish robustness to arbitrary or adversarial likelihood errors.

Second, a cause missing from the development data produces confident wrong verdicts. Its row is uninformative, so the observation that would reveal it counts as weak evidence for another cause: on sec-triage, four of eight causes are named as a wrong cause in every episode, and on sigma-triage 3 of 38. Every such verdict named a cause of the same class (a missing attack as another actionable cause, a missing benign explanation as another benign one), so no missing cause led to an attack being closed as benign in these tests.

Third, the v1.1 study adds errors that random mixing does not produce (sec-triage, 48 LLM-free episodes per cell; no attack closed as benign in any of its 3{,}968 episodes). \emph{Systematic confusion}, which mixes each cause's rows toward one fixed other cause, leaves the posterior stop almost unaffected at weight 0.25 (47 of 48) but blocks the confirmation stop entirely: no outcome remains ten times more likely under one cause than under its partner, so all 48 episodes spend the full budget and escalate. The confirmation rule thus keeps errors low by giving up automation when the table blurs two causes, and lowering its ratio to 3 does not recover it (0 of 48; \cref{tab:v12policy}), whereas joint evidence verifies 47 of 48 without a wrong verdict. \emph{A cause absent from the candidate set} is named as another actionable cause in 24 of 48 episodes under the posterior stop and escalated in all 48 under the confirmation stop. \emph{A duplicated source}, a second probe that reads the same feed as \texttt{source\_ips}, costs nothing measurable. Across thresholds from 0.6 to 0.95 the confirmation stop keeps 1.000 at rising cost (2.54 to 4.14 units), while the posterior stop falls to 0.875 from 0.9 on.

\subsection{Second Task Family}
\label{sec:external}
On sigma-triage, whose candidate causes come from the authors of public detection rules and whose outcomes were annotated by a model outside both planner families, both pre-registered endpoints on Qwen2.5-7B hold (98.33\% intervals over 12 rules; \cref{tab:v10a}, Appendix~\ref{app:tables}): the ranking adds $+0.171$ $[+0.094,+0.250]$, and \hesp{} with a controller stop exceeds Memory-only by $+0.296$ $[+0.088,+0.473]$. The ranking effect replicates on all five models ($+0.105$ to $+0.230$), and the separation of probing from stopping replicates: without a controller stop, Llama-3.1-8B finishes in one of 304 episodes, and a stop lifts it to 1.000. These configurations have no finish guard, and Qwen2.5-7B closed 10 of 144 attack episodes as benign inside \hesp{} (none under Memory-only), nine on one rule; in each, the ledger gave the claimed cause 0.002 and none of the cited observations supported it, so the guard would have rejected every one. The sec-triage primary endpoints are also robust to the clustering unit: resampling the 8 causes instead of the 24 tasks leaves every confirmed endpoint's interval above zero, and the effects are similar on static and drifting tasks (Appendix~\ref{app:tables}).

\subsection{Fit of Public Security Data to This Evaluation}
\label{sec:realdata}
Would these results hold on real data? We asked whether the three public datasets that come closest could test an investigation policy as defined here, which needs alerts that leave several explanations plausible, probes that differ in cost and informativeness, and objective ground truth. Appendix~\ref{app:public} gives versions, files, and methods; the ExCyTIn feasibility checks read its test questions (erratum E-10), whereas GUIDE's test split was never read. In the subsets we examined and with our adaptations, we found no setting that fits this question. In ExCyTIn-Bench~\citep{excytin}, most training questions name the alert to locate (54.8\%), and the one indicator-pivot probe we tested finds the answer no more often than distractors; this does not show that the benchmark as a whole tests retrieval only. In GUIDE~\citep{guide}, organization identity carries 1.03 of the grade's 1.50 bits, which may reflect policy, base rates, or detector mix, and evidence integration does not rank cold-start true positives better than history lookup (AUROC 0.755 against 0.759). In the 27 OTRF~\citep{otrf} remote-execution recordings and two event types we audited, a rule on the alert text separates 44 of 46 alerts. Paired attacks and legitimate administration that use the same tools, under objective labels, are what these subsets lack. The Cyber Defense Benchmark~\citep{cyberdefensebench} shows that such telemetry can be made interactive, as SQL threat hunting without a priming alert.
